\section{总结}
\label{sec:conclusion}

\subsection{可扩展之处}
资源优化优先针对运动放大和ROI模块开展缓存生存期分析，评估工作区复用、位宽调整和运算单元共享。在相同输入与约束下，同时比较数值误差、处理速率、BRAM/DSP占用和时序余量，确定各方案的收益与代价。

测量扩展按频率、纹理和振幅组织对照实验，引入参考传感器及成像比例标定，将像素位移转换为物理位移；在此基础上扩展可配置ROI与多区域分析。

下一阶段部署工作包括2025.2基线整机复测、冷启动验收、屏端源工程归档，以及HLS源码重生成与导出IP的等价比较，形成源码到板级运行的完整版本记录。

\subsection{心得体会}
\paragraph{软硬件协同与资源权衡}
视觉算法部署到板端后，计算、存储和任务调度相互影响。运动放大与ROI工作区是BRAM的主要开销来源，PS还要协调采集、计算和显示对帧缓存的使用。将模块资源分摊与帧完成记录结合分析，能够把资源开销对应到具体功能，并为缓存复用和任务调度调整提供依据。

\paragraph{大模型协作与证据核查}
在本次报告编制中，设计者确定应用定位、AMD评分重点、模板和素材；大模型辅助检索源码、整理数据、绘图及编写LaTeX。协作结果分别通过代码核对、原始数据检查、编译及PDF渲染确认。表\ref{tab:collaboration}列出三个具体案例，说明需求提出、修改执行与结果检查之间的对应关系。

\begin{table}[htbp]
  \centering\small
  \caption{报告编制中的协作案例与验证依据}
  \label{tab:collaboration}
  \begin{tabularx}{\linewidth}{@{}p{0.16\linewidth}Xp{0.22\linewidth}p{0.22\linewidth}@{}}
    \toprule
    实际问题 & 协作处理 & 检查依据 & 形成的方法 \\
    \midrule
    算法与版本标注 & 区分空间FFT、时间DFT、v18实板与2025.2实现 & 源码、数据及构建报告 & 每项指标记录算法、版本和条件 \\
    FPGA资源分析 & 汇总资源总量、模块分摊及BRAM余量 & 实现报告、模块资源表 & 从资源分布定位优化对象 \\
    图形与界面排版 & 横排流程图，更新串口屏UI并检查字号与分页 & 图形源文件、原始图像及PDF & 按最终插入尺寸检查可读性 \\
    \bottomrule
  \end{tabularx}
\end{table}

\paragraph{从工程经验到可复用技能包}
仓库将工程调试经验整理为三份指南\cite{repo}。《证据驱动的工程交接》从源码、产物哈希和测试目录建立版本对应关系；《串口控制与设备状态验证》依据原始字节、协议和状态表，逐层检查接收、解析、调度及完成事件；《启动镜像与冷启动验证》按镜像组成、传输、读回、自启动和外设业务划分检查步骤。

指南统一记录输入、操作步骤、完成判据和适用条件。文件核对、UART命令、实体触摸和外设运行分别保留对应结果，便于从异常现象追溯到具体环节。更换器件、工具链或协议时，先调整配置并核对适用条件，再执行相应步骤。

\Needspace{8\baselineskip}
\paragraph{版本一致性与复现经验}
源码、IP、初始化文件、位流和应用构成一套相互关联的交付物。复现文档同步记录工具版本、构建入口、输入参数、预期结果及原始日志路径；实板指标与实现报告分别注明版本和测试条件，使复现者能够沿同一入口检查功能、性能和产物对应关系。
